\section{Ba cái giá đọc cạnh nhau}
\label{sec:ch16-ba-cai-gia}

\index{lối thoát!ba cái giá không cùng hạng|(}%
Mục trước khép lại ở chỗ văn liệu của đường thứ ba không hỏi làm sao biết các
tên gán ra là đúng. Mục này đặt chỗ ấy cạnh hai chỗ tương ứng của hai đường kia,
và cái hiện ra khi đặt cạnh nhau là điều mà đọc rời từng đường không thấy được.

\index{lối thoát!ba đường trả lời ba câu}%
Ba cái giá không cùng một hạng, và đó là kết luận thứ nhất. Đường thứ nhất trả
lời đúng câu hỏi của định nghĩa~\ref{def:ch01-do-trung-thuc}, và trả lời trong
một phạm vi mà chính thiết kế đã cắt ra: bảng phân bố nhãn mà
mục~\ref{sec:ch16-gia-duong-dung-san} đọc cho 486 trong 1946 nhãn và một trong
bốn loại bước, còn ví dụ hai ngưỡng cho thấy phạm vi ấy hẹp ngay trên một hàm
hai biến. Đường thứ hai rẻ, chạy được, và trả lời một câu khác: nó xác nhận được lời
đề xuất rằng một chỉ số thay cho chất lượng mà người dùng cảm nhận, và không xác
nhận được lời đề xuất rằng một chỉ số đo độ trung thực. Đường thứ ba không trả
lời câu nào cho tới khi có ai kiểm định dụng cụ mà nó dùng để đọc tên. Một đường
thu hẹp câu hỏi, một đường đổi câu hỏi, một đường hoãn câu hỏi.

\index{lối thoát!ba đường không song song}%
Kết luận thứ hai thì mạnh hơn, và nó là kết luận mà nhan đề chương không gợi ra.
Ba đường không phải ba cửa song song, và ba bài báo của chúng nói vậy, mỗi bài
một câu.

Bài của đường thứ nhất chọn đường của nó \emph{vì} nó xét đường thứ ba là không
dùng được. Mục 3 của nó mở bằng đúng lập luận ấy: muốn biết mô hình đã tính gì
thì phải nhìn vào trong, các tính toán bên trong không quan sát trực tiếp được,
và các công cụ diễn giải hiện nay không cho phép truy cập một cách đáng tin cậy
và mở rộng được vào quá trình suy luận đứng sau một đầu ra; vì thế bài vòng qua
bằng thiết kế nhiệm vụ~\autocite{p21metrics}. Đường thứ nhất, trong bài đề xuất
ra nó, là hệ quả của việc đường thứ ba bị loại.

Bài của đường thứ ba nói chỗ thiếu dụng cụ đo là chỗ thiếu của cả ba nhánh
attribution cùng lúc, liệt kê ba cách đánh giá trong đó hai cách là hai thứ
chương này đang cân, rồi nói cả ba đều chưa đủ~\autocite{p18unified}. Nó không
nêu ground truth nào cho nhánh của đường thứ ba. Vậy đường ấy không đưa ai ra
khỏi phòng; nó là một bức tường khác của cùng căn phòng.

Bài của đường thứ hai bác cái giá mà nhan đề chương gán cho nó, và bác bằng lập
luận rằng những khó khăn ấy không phải là không vượt qua
nổi~\autocite{p24humangap}. Nghĩa là chỗ tắc của đường ấy không nằm ở chỗ nó
đắt.

\index{lối thoát!bài nào cũng chỉ vào bài khác}%
Ba câu ấy ghép lại thành một hình mà không bài nào tự vẽ. Người đề xuất đường
thứ nhất chỉ sang đường thứ ba và nói không đi được. Người vẽ bản đồ đường thứ
ba nói chỗ thiếu nằm ở cả ba nhánh. Người đếm đường thứ hai nói đường ấy đi
được nhưng gần như không ai đi, và không phải vì đắt. Ba lối thoát không xếp
thành ba cửa cho một người chọn; chúng xếp thành một vòng, trong đó mỗi bên chỉ
sang bên khác. Câu ấy là câu của sách: không bài nào trong bốn bài phát biểu nó,
vì mỗi bài chỉ đứng ở một chỗ trong vòng.

\index{lối thoát!chữ cái giá đọc lại}%
Kết luận thứ ba là về chữ. Chữ \enquote{cái giá} gợi ra một thứ đo bằng tiền,
bằng thời gian, hoặc bằng lượt chạy, và trong ba đường thì không đường nào trả
giá kiểu ấy. Đường thứ nhất mất tính đại diện. Đường thứ hai mất chính câu hỏi.
Đường thứ ba mất chỗ đứng, vì dụng cụ nó dùng để đọc lại là thứ chưa ai kiểm.
Ba khoản mất ấy đều là khoản mất về nhận thức chứ không về tài nguyên, và đó là
lý do không đường nào rẻ đi khi máy tính nhanh hơn. Đọc theo
bảng~\ref{tab:ch01-bon-quan-he}, cả ba đều là chuyện dòng nào của bảng được cấp
bằng chứng, chứ không phải chuyện bằng chứng ấy tốn bao nhiêu.

\index{khoảng trống nghiên cứu!cái chương này để lại}%
Ba thứ đi tiếp từ chương này, và chúng đi tới hai chỗ khác nhau.

Thứ đi tiếp thứ nhất đi tới chương~\ref{ch:khoang-trong-nghien-cuu}, và nó là cái vòng
vừa mô tả cùng với ba khoản mất của nó. Chương ấy lấy giao của các phát biểu hạn
chế, và một lối thoát mà mỗi bên chỉ sang bên khác là một dữ kiện về hình dạng
của khoảng trống chứ không phải một chỗ trống trong sổ.

Thứ đi tiếp thứ hai cũng đi tới đó, và nó là yêu cầu thứ ba mà
mục~\ref{sec:ch15-doc-duoc-dieu-khien-duoc} đã gom, nay đọc lại được ở mức
chung. Ba yêu cầu ấy, mô tả chỗ lệch từ mục~\ref{sec:ch13-tran-noi-gi}, công bố
tập can thiệp từ mục~\ref{sec:ch14-dieu-chua-giai-quyet}, và báo cả phép đọc tên
lẫn phép ép từ mục~\ref{sec:ch15-doc-duoc-dieu-khien-duoc}, đều là yêu cầu đặt
lên một dụng cụ tương lai. Chương này thêm điều kiện thứ tư và nó đến từ chính
phép rà ở mục~\ref{sec:ch16-duong-thu-ba}: một dụng cụ đo dựng trên một dụng cụ
khác phải nói dụng cụ kia đã được kiểm định bằng gì. Đó là chỗ mà cả đường thứ
ba lẫn phép sửa của nó cùng hụt, và là chỗ mà
chương~\ref{ch:tro-choi-attribution} đã gặp một vòng cùng hình dạng ở một tầng
khác.

\index{lối thoát!ba cái giá không cùng hạng|)}%
\index{chuỗi suy luận!lối thoát thứ tư}%
Thứ đi tiếp thứ ba đi tới chương sau, và nó là một lối thoát mà chương này cố tình không
đọc. Luận đề ở mục~\ref{sec:ch01-luan-de-cua-sach} kể bốn đường mà văn liệu phê
phán chỉ tới, không phải ba; đường thứ tư là chuyển cả câu hỏi sang
\tn{chuỗi suy luận}{chain of thought}, tức sang những mô hình tự viết ra lời
giải thích cho chính mình bằng ngôn ngữ. Đường ấy khác ba đường trên ở một chỗ
mà cả chương này chưa chạm tới: ở ba đường trên, lời giải thích do một công cụ
bên ngoài dựng, còn ở đường thứ tư thì chính mô hình sản xuất nó. Điều đáng chú
ý là bài của đường thứ nhất đã ở sẵn bên kia ranh giới ấy, vì đối tượng nó đo là
chuỗi suy luận chứ không phải attribution đặc trưng, và
mục~\ref{sec:ch09-doi-tuong-khac-nhau} đã giữ ranh giới ấy suốt bảy chương.
Chương~\ref{ch:do-trung-thuc-chuoi-suy-luan} bước qua nó, và câu hỏi nó mang
theo từ đây là câu hỏi mà ba đường vừa đọc đều không đặt ra được: khi lời giải
thích do chính mô hình viết ra, thì cái được so với hành vi là ai nói.
